\documentclass[11pt,oneside]{article}

\usepackage[a4paper,margin=27mm,headheight=15pt]{geometry}
\usepackage[T1]{fontenc}
\usepackage[utf8]{inputenc}
\IfFileExists{libertinus.sty}{\usepackage{libertinus}}{\usepackage{lmodern}}
\usepackage{microtype}
\usepackage{amsmath,amssymb,amsthm,mathtools,mathrsfs}
\usepackage{bm}
\usepackage{booktabs,longtable,array,tabularx}
\usepackage{graphicx}
\usepackage{pgfplots}
\usepackage{xcolor}
\usepackage{enumitem}
\usepackage{fancyhdr}
\usepackage{titlesec}
\usepackage{listings}
\usepackage{xurl}
\usepackage{hyperref}
\usepackage[nameinlink,noabbrev]{cleveref}

\definecolor{linkblue}{RGB}{25,75,135}
\definecolor{shadegray}{RGB}{246,247,249}
\definecolor{rulegray}{RGB}{195,201,208}
\pgfplotsset{compat=1.18}
\hypersetup{
  colorlinks=true,
  linkcolor=linkblue,
  citecolor=linkblue,
  urlcolor=linkblue,
  pdftitle={Primary Exposition Gap Register},
  pdfsubject={Unformalized claims removed from the Fabius primary exposition},
  pdfkeywords={Fabius function, research frontier, formalization gap, Lean}
}

\pagestyle{fancy}
\fancyhf{}
\fancyhead[L]{\small Fabius--Rvachev research frontier}
\fancyhead[R]{\small Primary exposition gap register}
\fancyfoot[C]{\thepage}
\renewcommand{\headrulewidth}{0.4pt}

\titleformat{\section}{\Large\bfseries}{\thesection}{0.7em}{}
\titleformat{\subsection}{\large\bfseries}{\thesubsection}{0.7em}{}
\titleformat{\subsubsection}{\normalsize\bfseries}{\thesubsubsection}{0.7em}{}
\setcounter{secnumdepth}{3}
\setcounter{tocdepth}{2}
\setlist[itemize]{leftmargin=2em,itemsep=0.25em,topsep=0.4em}
\setlist[enumerate]{leftmargin=2.3em,itemsep=0.25em,topsep=0.4em}

% Long unbreakable inline mathematics occasionally leaves a line with no legal
% break point.  A small emergency stretch lets TeX loosen such a line rather
% than let it protrude into the margin; it applies only to lines that would
% otherwise be overfull.
\setlength{\emergencystretch}{2em}

\newtheorem{theorem}{Theorem}[section]
\newtheorem{proposition}[theorem]{Proposition}
\newtheorem{lemma}[theorem]{Lemma}
\newtheorem{corollary}[theorem]{Corollary}
\newtheorem{conjecture}[theorem]{Conjecture}
\theoremstyle{definition}
\newtheorem{definition}[theorem]{Definition}
\newtheorem{algorithm}[theorem]{Algorithm}
\newtheorem{example}[theorem]{Example}
\theoremstyle{remark}
\newtheorem{remark}[theorem]{Remark}
\newtheorem{warning}[theorem]{Warning}

\newcommand{\R}{\mathbb R}
\newcommand{\C}{\mathbb C}
\newcommand{\Q}{\mathbb Q}
\newcommand{\N}{\mathbb N}
\newcommand{\Z}{\mathbb Z}
\newcommand{\Up}{\operatorname{up}}
\newcommand{\sinc}{\operatorname{sinc}}
\newcommand{\wt}{\operatorname{wt}_2}
\newcommand{\e}{\mathrm e}
\newcommand{\ii}{\mathrm i}
\newcommand{\dd}{\,\mathrm d}
\newcommand{\Li}{\operatorname{Li}}
\newcommand{\Var}{\operatorname{Var}}
\newcommand{\E}{\mathbb E}
\newcommand{\Prob}{\mathbb P}
\newcommand{\bigO}{\mathcal O}
\newcommand{\smallo}{o}
\newcommand{\supp}{\operatorname{supp}}
\newcommand{\sgn}{\operatorname{sgn}}
\newcommand{\lcm}{\operatorname{lcm}}
\newcommand{\vtwo}{v_2}
\newcommand{\qbinom}[3]{\genfrac{[}{]}{0pt}{}{#1}{#2}_{#3}}
\newcommand{\repo}{\nolinkurl{Analysis/FabiusFunction}}
\newcommand{\lean}[1]{%
  \ifmmode\text{\nolinkurl{#1}}\else\nolinkurl{#1}\fi}

\newenvironment{proofidea}{\par\smallskip\noindent\textit{Proof architecture.}\ }{\hfill$\square$\par\smallskip}
\newenvironment{boxedremark}{%
  \par\smallskip\noindent\begin{center}\begin{minipage}{0.94\textwidth}%
  \color{black}\hrule\vspace{0.6em}\small
}{%
  \vspace{0.6em}\hrule\end{minipage}\end{center}\smallskip
}

\lstdefinestyle{pseudo}{
  basicstyle=\ttfamily\small,
  backgroundcolor=\color{shadegray},
  frame=single,
  rulecolor=\color{rulegray},
  columns=fullflexible,
  keepspaces=true,
  showstringspaces=false,
  xleftmargin=0.7em,
  xrightmargin=0.7em,
  aboveskip=0.8em,
  belowskip=0.8em
}

% Frontier-specific environments extend, rather than replace, the primary
% document's theorem environment set.
\newtheorem{candidate}[theorem]{Candidate statement}
\newtheorem{obligation}[theorem]{Formalization obligation}

\title{\bfseries Primary Exposition Gap Register\\[0.35em]
\large Claims removed pending exact Lean counterparts}
\author{A research-frontier companion to
\emph{The Fabius Function and Rvachev's Up-Function}\\
\small based on Vladimir Reshetnikov's \texttt{ProveIt} formal development}
\date{Repository snapshot: 25 August 2026}

\begin{document}
\maketitle
\begin{abstract}
This notebook preserves mathematical material removed from the primary Fabius exposition
during its proof-correspondence audit.  Nothing here is claimed as a theorem of the Lean
development unless explicitly described as formal background.  Each item needs a new
public theorem, a precise cost or convergence model, or a correction to its statement.
The larger endpoint-asymptotic frontier is maintained separately in the sibling notebook
\emph{Small-Argument Asymptotics}.
\end{abstract}
\tableofcontents
\newpage

\section{Status and promotion rule}

The primary exposition admits a mathematical statement only when an actual proved Lean
declaration has the same hypotheses and conclusion.  An elementary consequence, a paper
proof, a source comment, an executable calculation, or an equivalent normalization is not
sufficient.  Promotion from this register requires a public Lean theorem, a successful
build, an exact source citation in the exposition, and a rebuilt, visually checked PDF.

\section{Fourier-decay normalization}

The source proves polynomially weighted decay.  The conventional shifted-weight form that
was removed is:
\begin{candidate}
For every \(N\in\N\), there is \(C_N>0\) such that
\[
 |\widehat{\Up}(\xi)|\le C_N(1+|\xi|)^{-N}\qquad(\xi\in\R).
\]
\end{candidate}
It is standardly equivalent to the proved estimate after a bounded-region argument and a
change of constant, but it has no literal named Lean bridge.
\begin{obligation}
Derive the shifted-weight estimate, including \(|\xi|\le1\), and expose it as a public
declaration.
\end{obligation}

\section{Dyadic specializations and evaluator cost}

\subsection{Concrete rational values}

Let \(F\) denote the bounded Fabius function and \(\Up\) Rvachev's compactly supported
up-function.  The notation
\(c_n=\int_{-1}^{1}x^{2n}\Up(x)\dd x\) is the even-moment notation of the primary article.
The general dyadic evaluator is proved correct, but these specializations have no named
theorems:
\begin{candidate}
\[
 c_1=\frac19,\qquad
 F(1/4)=\frac5{72},\qquad
 F(1/8)=\frac1{288},\qquad
 F(3/8)=\frac{73}{288}.
\]
\end{candidate}
\begin{obligation}
Prove named specializations by reducing the exact evaluator and bridge each result to the
analytic function.
\end{obligation}

\subsection{Popcount complexity}

Termination is formalized, but the removed prose did not distinguish the following two
possible complexity claims.  For \(e,a\in\N\) with \(0<a<2^e\), let \(D(e,a)\) be the
number of recursive calls that enter the evaluator's interior branch.
\begin{candidate}
The exact-count candidate is \(D(e,a)=\wt(a)\).
\end{candidate}
\begin{candidate}
The weaker upper-bound candidate is \(D(e,a)\le\wt(a)\).
\end{candidate}
\begin{obligation}
Define \(D\) in Lean, settle which statement matches the executable recursion, and prove
the resulting relation to binary weight, including endpoint and saturation branches.
\end{obligation}

\subsection{Independent analytic derivations}

Earlier prose derived the dyadic formula through Taylor integrals, block substitutions,
and the exact remainder \(-F(y)\).  The final rational identities are formal, but those
intermediate analytic equations are not public theorem-for-theorem counterparts.
\begin{obligation}
Formalize named Taylor integral identities for the bounded and signed global functions,
including the blockwise changes of variables and the exact remainder.
\end{obligation}

\section{Complex-shift spline estimate}

The proved estimate has the factor
\(\exp(|Q-\tfrac12|)\).  Put
\[
 N_p(x)=\max\!\left(0,\left\lfloor2^px+\frac12\right\rfloor\right),\qquad
 S_p^{(Q)}(x)=\frac{(-1)^p}{2^{\binom p2}p!}
 \sum_{r=0}^{N_p(x)-1}(-1)^{\wt(r)}(r-2^px+Q)^p,
\]
with an empty sum when \(N_p(x)=0\).  The removed proposed sharpening was:
\begin{candidate}
For every \(p\in\N\) with \(p\ge1\), every \(Q\in\C\), and every \(x\in\R\),
\[
 \left|S_p^{(Q)}(x)-S_p^{(1/2)}(x)\right|
 \le2^{-(p-1)}
 \left(\exp\!\left(\left|Q-\frac12\right|\right)-1\right).
\]
\end{candidate}
\begin{warning}
The accompanying assertion that the relevant triangular exponent attains equality
\emph{exactly in the top term} is false when \(p=2\): both endpoint indices attain the same
exponent.
\end{warning}
\begin{obligation}
Formalize the finite Taylor majorant, state the minimizer edge cases correctly, and prove a
public \(\exp-1\) estimate if it survives that correction.
\end{obligation}

\section{Prefix-kernel interpretations}

The binomial kernel and the resulting prefix-sum formula are formalized.  These two
interpretations are not:
\begin{candidate}
\begin{itemize}
\item for \(i,k\in\N\) with \(k\ge1\), \(\binom{i+k-1}{k-1}\) counts ordered \(k\)-tuples
      of nonnegative integers summing to \(i\);
\item for \(k\in\N\) with \(k\ge1\), the kernel is the \(k\)-fold discrete convolution
      power of the constant kernel.
\end{itemize}
\end{candidate}
\begin{obligation}
Choose a finite-support or locally finite convolution model, prove the counting bijection,
and connect it to the existing \texttt{iteratedPrefixKernel}.
\end{obligation}

\section{Block-polynomial derivatives at one}

The finite factorization and its power-sum consequences are formal.  The literal derivative
version removed from the article was:
\begin{candidate}
For \(m,r\in\N\) with \(m<r\), the \(m\)-th derivative of
\(\prod_{j<r}(1-z^{2^j})\) vanishes at \(z=1\), while its \(r\)-th derivative there is
\[
 (-1)^r r!\,2^{\binom r2}.
\]
\end{candidate}
\begin{obligation}
Formalize polynomial iterated derivatives at \(z=1\) and identify them with the proved
coefficient or power-sum formula.
\end{obligation}

\section{An actual infinite formal product}

Lean proves coefficient stabilization of the finite products
\(\prod_{j<r}(1-z^{2^j})\).  The stronger construction formerly asserted was:
\begin{candidate}
Construct \(\prod_{j\ge0}(1-z^{2^j})\) as a formal infinite product and prove that it equals
the Thue--Morse power series.
\end{candidate}
\begin{obligation}
Select an appropriate formal-power-series topology or summability API, define the product,
and prove equality with the coefficient-stabilized series.  Until then, the notation is
only shorthand for finite stabilization.
\end{obligation}

\section{Formal-background crosswalk}

The following declarations support the proved background cited above.  None of them is a
counterpart for the adjacent candidate statement.

\begingroup
\small
\setlength{\tabcolsep}{0.6em}
\noindent\begin{tabularx}{\textwidth}{@{}>{\raggedright\arraybackslash}p{0.29\textwidth}
  >{\raggedright\arraybackslash}X@{}}
\toprule
Background & Lean declaration and module\\
\midrule
Polynomially weighted Fourier decay &
\lean{rvachevFourier_real_rapidDecay} in \path{PoissonSummation.lean}\\
Even moments and exact dyadic evaluation &
\lean{moment}, \lean{fabiusDyadicValue}, \lean{fabiusDyadicValue_cast},
\lean{fabiusDyadic_cast} in \path{Arithmetic.lean}, \path{PaperStatements.lean}, and
\path{DyadicAnalytic.lean}\\
Complex-shift exponential bound &
\lean{norm_fabiusComplexShiftSpline_sub_center_le_half_pow_mul_exp} in
\path{FabiusComplexShiftSpline.lean}\\
Prefix kernel and convolution formula &
\lean{iteratedPrefixKernel} and \lean{iteratedPrefix_eq_sum_kernel} in
\path{ThueMorsePrefix.lean}\\
Even/odd Thue--Morse recurrences &
\lean{binaryWeight_two_mul} and \lean{binaryWeight_two_mul_add_one} in
\path{DyadicClosedForm.lean}\\
Finite block product and coefficient stabilization &
\lean{thueMorseBlockPolynomial_eq_product} and \lean{coeff_finite_thueMorse_product} in
\path{ThueMorseGenerating.lean}\\
\bottomrule
\end{tabularx}
\endgroup

\section{Endpoint material maintained elsewhere}

The sibling notebook \emph{Small-Argument Asymptotics} preserves the removed raw
Mellin-inversion derivation, exponential Gamma--zeta decay estimates, numerical periodic
profiles, proposed ordered-composition formulas for saddle coefficients, explicit higher
Lambert polynomials, a proposed leading-derivative law, uniqueness and phase-density
arguments, half-moment asymptotics, and strengthened endpoint comparisons.

The primary article now retains only the exact negative-Laplace decomposition, the proved
Gamma--zeta Fourier series, the evaluated constant, lower-Lambert reduction, Bromwich
inversion, the first two periodic saddle corrections, and the recursive all-orders theorem.

\section{Provenance}

This register was created from the claim-by-claim audit of the former drafts and primary
exposition on 25 August 2026.  Related numerical and exploratory material is also preserved
in the dyadic, \(q\)-binomial, inverse, integration, and repeated-integration notebooks in
this research-frontier directory.

\end{document}
